% ===================== CHAPITRE 7 — SPRINT 5 =========================
\chapter{Sprint 5 : Chaîne ETL et volet décisionnel}

\section*{Introduction}
Ce sprint réalise le \textbf{volet décisionnel} de la solution. Il met en place
une chaîne \textbf{ETL} qui consolide les données opérationnelles dans un
\textbf{datamart}, puis restitue les indicateurs au moyen d'un tableau de bord
embarqué et de rapports Power BI.

\section{Objectifs du Sprint}
L'objectif est de \textbf{consolider} les données dispersées (employés, clients,
objectifs, pointage) en un datamart agrégé par agence, et de \textbf{restituer}
les indicateurs de pilotage aux responsables.

\section{Spécification des besoins}

\subsection{Diagramme de cas d'utilisation du Sprint 5}
La figure~\ref{fig:uc_s5} présente les cas d'utilisation du sprint.

\begin{figure}[htbp]
  \centering
  \includegraphics[width=0.86\textwidth]{uc_sprint5.png}
  \caption{Diagramme de cas d'utilisation du Sprint 5}
  \label{fig:uc_s5}
\end{figure}

\subsection{Sprint Backlog du sprint 5}
Le tableau~\ref{tab:bl_s5} détaille les user stories et les tâches techniques.

\begin{table}[htbp]
\centering
\caption{Sprint Backlog du sprint 5}
\label{tab:bl_s5}
\renewcommand{\arraystretch}{1.3}
\begin{tabularx}{\textwidth}{>{\bfseries}p{1cm}p{5cm}X}
\toprule
\textmd{\textbf{US}} & \textbf{User story} & \textbf{Tâches techniques}\\
\midrule
US14 & Alimenter le datamart via la chaîne ETL. &
Développer le script Python \texttt{etl\_agences.py} (pandas) ; étapes
\emph{Extract / Transform / Load} ; agréger les indicateurs par agence ; produire
le datamart en étoile (\texttt{dim\_agence}, \texttt{fait\_agence}).\\
\midrule
US15 & Consulter le tableau de bord décisionnel. &
Enrichir le \texttt{HomeServlet} avec les statistiques ; réaliser le tableau de
bord \texttt{home.jsp} (cartes d'indicateurs et graphiques ApexCharts).\\
\midrule
US16 & Consulter les rapports Power BI. &
Créer les vues analytiques \texttt{V\_BI\_EMPLOYES}, \texttt{V\_BI\_CLIENTS},
\texttt{V\_BI\_OBJECTIFS}, \texttt{V\_BI\_POINTAGE} ; les connecter à Power BI.\\
\bottomrule
\end{tabularx}
\end{table}

\section{Analyse des besoins}

\subsection{Description textuelle des cas d'utilisation}

\cudesc{tab:cu_etl}{Alimenter le datamart (ETL)}
{Consolider les données opérationnelles en un datamart agrégé par agence.}
{Processus ETL / Administrateur}
{Les données sources (agences, employés, clients, objectifs, pointage) doivent
être disponibles.}
{Le datamart en étoile est produit et mis à la disposition de la restitution et de
la segmentation.}
{(1) L'étape \emph{Extract} lit les données sources ; (2) l'étape \emph{Transform}
les nettoie et les agrège par agence ; (3) l'étape \emph{Load} écrit le datamart
(\texttt{dim\_agence}, \texttt{fait\_agence}).}

\cudesc{tab:cu_dashboard}{Consulter le tableau de bord}
{Restituer aux responsables une vision synthétique de l'activité.}
{Administrateur / Directeur commercial}
{L'acteur doit être authentifié.}
{Les indicateurs clés (effectifs, présence, production) sont affichés.}
{(1) L'acteur se connecte ; (2) le \texttt{HomeServlet} calcule les statistiques ;
(3) la page \texttt{home.jsp} affiche les cartes d'indicateurs et les graphiques
ApexCharts.}

\section{Conception}

\subsection{Chaîne décisionnelle et processus ETL}
Les données nécessaires au pilotage sont dispersées dans les tables opérationnelles
et exprimées à des grains hétérogènes. Le processus \textbf{ETL} (\emph{Extract --
Transform -- Load}) les consolide en un datamart agrégé par agence. La
figure~\ref{fig:etl_chaine} présente cette chaîne décisionnelle.

\begin{figure}[htbp]
  \centering
  \includegraphics[width=\textwidth]{etl_chaine.png}
  \caption{Chaîne décisionnelle : du système opérationnel au datamart et à sa restitution}
  \label{fig:etl_chaine}
\end{figure}

Trois étapes s'enchaînent : \textbf{Extract} (lecture des sources
\texttt{AGENCE}, \texttt{B\_UTILISATEURS}, \texttt{CLIENT}, \texttt{B\_OBJECTIF},
\texttt{POINTAGE}), \textbf{Transform} (nettoyage, typage et agrégation par agence)
et \textbf{Load} (chargement du datamart en étoile). Ce datamart constitue une
\textbf{source unique consolidée}, partagée par la restitution (tableau de bord et
Power BI) et par la segmentation des agences (chapitre~\ref{chap:segmentation}).

\subsection{Modèle en étoile de l'entrepôt}
L'entrepôt est organisé selon une logique en \textbf{étoile}
(figure~\ref{fig:bi}) : autour de tables de faits (production, effectifs, clients,
présence), des dimensions descriptives (agence, gestionnaire, période, type de
client) permettent l'analyse multidimensionnelle. Les vues \texttt{V\_BI\_*}
matérialisent cette couche sémantique pour l'outil de restitution.

\begin{figure}[htbp]
  \centering
  \includegraphics[width=0.95\textwidth]{bi.png}
  \caption{Modèle décisionnel en étoile (vues \texttt{V\_BI\_*})}
  \label{fig:bi}
\end{figure}

Le tableau~\ref{tab:kpi} récapitule les principaux indicateurs (KPI) retenus.

\begin{table}[htbp]
\centering
\caption{Indicateurs de performance (KPI) du volet décisionnel}
\label{tab:kpi}
\renewcommand{\arraystretch}{1.25}
\begin{tabularx}{\textwidth}{>{\bfseries}p{4.2cm}Xp{2.6cm}}
\toprule
\textmd{\textbf{Indicateur}} & \textbf{Description} & \textbf{Source}\\
\midrule
Effectifs par agence & Nombre d'employés par agence & \texttt{V\_BI\_EMPLOYES}\\
Répartition des rôles & Gestionnaires / conseillers / hors commercial & \texttt{V\_BI\_EMPLOYES}\\
Portefeuille clients & Clients par type et par agence & \texttt{V\_BI\_CLIENTS}\\
Ouvertures de comptes & Comptes chèques + épargne + courants & \texttt{V\_BI\_OBJECTIFS}\\
Production de crédits & Conso, immobilier, investissement & \texttt{V\_BI\_OBJECTIFS}\\
Collecte d'épargne & Épargne additionnelle collectée & \texttt{V\_BI\_OBJECTIFS}\\
Taux de présence & Présents, retards et absences par agence et période & \texttt{V\_BI\_POINTAGE}\\
\bottomrule
\end{tabularx}
\end{table}

\section{Réalisation}

\subsection{Mise en œuvre de la chaîne ETL en Python}
La chaîne ETL est implémentée en \textbf{Python} (script \texttt{etl\_agences.py})
à l'aide de la bibliothèque \textbf{pandas}~\cite{pandas}. Elle est organisée en
trois fonctions : \texttt{extract()} lit les données sources (connexion Oracle via
le connecteur \texttt{oracledb} ou, à défaut, export CSV) ; \texttt{transform()}
nettoie et \textbf{agrège les données par agence} pour produire les sept
indicateurs de pilotage ; \texttt{load()} écrit le datamart en étoile
(\texttt{dim\_agence}, \texttt{fait\_agence}) et le fichier réutilisé par la
segmentation. L'extrait~\ref{lst:etl_transform} illustre le cœur de l'agrégation.

\begin{lstlisting}[style=pycode,caption={Agrégation des indicateurs par agence (étape Transform)},label={lst:etl_transform}]
# effectif et nombre de gestionnaires par agence
eff = employes.groupby("SK_AGENCE").agg(
    effectif=("SK_UTILISATEUR", "count"),
    nb_gestionnaires=("EST_GESTIONNAIRE", "sum")).reset_index()
# portefeuille clients
cli = clients.groupby("SK_AGENCE").size().reset_index(name="nb_clients")
# production (comptes, credits, epargne)
obj = objectifs.groupby("SK_AGENCE").agg(
    total_comptes=("COMPTES", "sum"),
    production_credits=("CREDITS", "sum"),
    collecte_epargne=("EPARGNE", "sum")).reset_index()
# taux de presence issu du pointage
poi = pointages.assign(
        present=pointages["STATUT"].isin(["PRESENT", "RETARD"]).astype(int)) \
    .groupby("SK_AGENCE").agg(taux_presence=("present", "mean")).reset_index()
\end{lstlisting}

L'exécution du script produit le datamart sous \texttt{etl/entrepot/}
(\texttt{dim\_agence.csv}, \texttt{fait\_agence.csv}) ainsi que le fichier
consolidé \texttt{clustering/data/agences.csv}. Le pipeline est \textbf{rejouable}
et journalise le nombre d'enregistrements traités à chaque étape.

\subsection{Tableau de bord et restitution Power BI}
Le tableau de bord embarqué (\texttt{home.jsp}, alimenté par le
\texttt{HomeServlet}) comporte une rangée de \textbf{cartes d'indicateurs}
(employés, clients, présents du jour, taux de présence, demandes en attente) puis
des \textbf{graphiques} ApexCharts (effectif par agence, répartition des rôles,
présence du jour). En complément, les vues \texttt{V\_BI\_*} sont connectées à
\textbf{Power BI} pour des rapports interactifs. La figure~\ref{fig:ui_dashboard}
en présente les interfaces. \emph{Les captures d'écran réelles seront insérées à
cet emplacement.}

\begin{figure}[htbp]
  \centering
  \fbox{\parbox[c][3.6cm][c]{0.8\textwidth}{\centering\itshape Capture à insérer
  -- tableau de bord et rapport Power BI}}
  \caption{Tableau de bord embarqué et rapport Power BI}
  \label{fig:ui_dashboard}
\end{figure}

\section*{Conclusion}
Ce sprint a doté la solution d'un volet décisionnel complet : une chaîne ETL
alimente un datamart consolidé, restitué par un tableau de bord et par Power BI.
Le datamart ainsi produit constitue l'entrée du dernier sprint, consacré à la
segmentation des agences par apprentissage automatique.
